\paragraph{Support labels.}
\begin{description}
  \item[\textbf{[Established]}] Follows from cited prior work or a complete
        argument given here.
  \item[\textbf{[Assumption-dependent]}] True under additional assumptions not
        derived from \textrm{OPH} first principles.
  \item[\textbf{[Conjecture]}] A plausible open direction, not a settled result.
\end{description}

% ------------------------------------------------------------
\subsection*{B.1\quad Theorem 1: QBFT Safety Bound}
% ------------------------------------------------------------

\begin{definition}[QBFT-style protocol]
\label{def:qbft-style}
A consensus protocol is \emph{QBFT-style} in this analysis if it satisfies
the following three structural properties.
The safety proof of Theorem~\ref{thm:qbft-safety} uses all three;
the theorem does not hold for protocols lacking any of them without a
compensating change to the argument.
\begin{enumerate}[(P1)]
  \item \textbf{One-vote-per-view.}
        Each nonfaulty node casts at most one vote per view number.
        A node that has voted in view~$v$ ignores any later request
        to vote in view~$v$.
  \item \textbf{Certificate semantics.}
        A decision requires a valid quorum certificate: in the classical
        exact-size case used below, $q=2f+1$ distinct, unforgeable,
        authenticated votes for the same value in the same view.  For
        larger validator sets at fixed fault budget, the quorum threshold must
        be scaled so that the overlap condition in (A6) remains true.
  \item \textbf{DLS-style view-change.}
        If no certificate is produced within a timeout, every nonfaulty node
        increments the view number by one and a new leader is selected by a
        fixed deterministic rule.  At \textrm{GST}, timeouts fire correctly
        and the view-change terminates in bounded rounds.
\end{enumerate}
The Istanbul BFT / QBFT protocol family
\cite{Moniz2020,Saltini}
satisfies (P1)--(P3) and is the intended instance.
\end{definition}

\paragraph{Assumptions A1--A6.}
\begin{enumerate}[(A1)]
  \item \textbf{Partial synchrony (DLS).}
        Fixed but initially unknown bounds $\Delta$ (message delay) and
        $\Phi$ (processing rates).
        \emph{Safety} holds without extra timing assumptions;
        \emph{liveness} holds after the Global Stabilisation Time (\textrm{GST}).
  \item \textbf{Byzantine fault model.}
        At most $f$ observers behave arbitrarily; the remaining $n-f$ are nonfaulty.
  \item \textbf{Classical exact sizing for this fixed-quorum theorem.}
        $n = 3f+1$, so that $q=2f+1$ certificates have a nonfaulty overlap
        witness.  The usual resilience condition $n\ge 3f+1$ is necessary for
        the fault model, but when $n>3f+1$ a fixed $q=2f+1$ quorum is
        insufficient for the overlap used in the safety proof; one must either
        impose (A6) directly or choose a threshold $q$ with
        $2q\ge n+f+1$.
  \item \textbf{Strong quorum connectivity.}
        Every quorum $Q$ with $|Q|=q$ is strongly connected within~$G$:
        for any $u,v\in Q$ there is a directed path in~$G$ contained
        entirely in~$Q$.
        This is strictly stronger than requiring the overlap graph of
        quorums to be connected, and is needed to propagate signed votes
        within a quorum.
  \item \textbf{Message authentication.}
        All messages carry unforgeable digital signatures.
  \item \textbf{OPH quorum overlap.}
        Any two quorums $Q_a, Q_b$ of size $q$ satisfy
        $|Q_a\cap Q_b|\geq f+1$.  For $q=2f+1$ this is guaranteed by the
        exact sizing $n=3f+1$ in (A3); for general $n$ it is the separate
        threshold condition $2q\ge n+f+1$, not a consequence of
        $n\ge3f+1$ alone.
\end{enumerate}

D.\ Matscheko's Lean audit of this appendix checks the finite quorum-overlap
core and records the same boundary caveat: at fixed \(q=2f+1\), the overlap
step is exact-size \(n=3f+1\) logic unless (A6) is imposed separately.

\begin{theorem}[QBFT Safety Bound
  {\normalfont\small [Same-view case established under A1--A6;
  cross-view safety conditional on a (P4) locking rule]}]
\label{thm:qbft-safety}
Under assumptions (A1)--(A6), any consensus protocol satisfying
(P1)--(P3) of Definition~\ref{def:qbft-style} and run over the
\textrm{OPH} observer graph satisfies:
\begin{enumerate}[(i)]
  \item \textbf{Same-view safety.} No two nonfaulty observers finalise
        conflicting patch states in the same view.
  \item \textbf{Liveness.} After \textrm{GST}, every nonfaulty observer
        finalises within $O(f\cdot\Delta)$ wall-clock time.
  \item \textbf{Optimality.} The bound $f<n/3$ is tight.
\end{enumerate}
\end{theorem}

\paragraph{Cross-view boundary.}
Clause (i) is stated for same-view finalisation because the argument below
covers exactly that case. Cross-view safety in the IBFT/QBFT protocol family
uses an additional prepared-certificate locking property, an explicit
(P4) hypothesis under which a nonfaulty node votes in a later view only for
a value carried by the highest prepared certificate it has received.
(P1)--(P3) alone do not exclude finalisation of conflicting values across
views; a cross-view extension of this theorem requires (P4) and its
accompanying view-change justification argument.
On its stated assumptions, this finalisation theorem is a finite carrier of stack row C1 of
the program's consistency stack: record permanence, meaning no two nonfaulty observers finalise
conflicting patch states, enters as a consistency requirement on the observer net, not as an
added postulate.

\begin{proof}[Proof sketch]
\textit{Safety.}
Suppose $O_a$ and $O_b$ finalise $s_a\neq s_b$ in the same view.
By (P2), each required a certificate of $q$ votes: sets $Q_a,Q_b$.
By (A6), $|Q_a\cap Q_b|\geq f+1$.  In the classical exact-size case this is
the inclusion-exclusion calculation
$|Q_a\cap Q_b|\geq(2f+1)+(2f+1)-(3f+1)=f+1$.
By (A2), at most $f$ are Byzantine, so $Q_a\cap Q_b$ contains a
nonfaulty $O^*$.
By (A4), $O^*$'s signed vote is path-reachable within both quorums.
By (P1), $O^*$ voted for at most one value. Contradiction.

\textit{Liveness and Optimality}
follow from \cite{DLS88} (Thm.~4.4) and \cite{Lamport1982}, cited directly.

\textit{Note on FLP.}
Fischer, Lynch, Paterson~\cite{FLP85} is an impossibility result for
fully asynchronous systems; it does not bear on achievability under
partial synchrony (A1).
\end{proof}

% ------------------------------------------------------------
\subsection*{B.2\quad Theorem 2: Convergence of the OPH Repair Map}
% ------------------------------------------------------------

\begin{definition}[OPH Repair Map: Petz form]
\label{def:oph-repair-map}
Let $\sigma\in\mathcal{D}(\mathcal{H})$ be a full-rank reference state and
$\mathcal{N}:\mathcal{B}(\mathcal{H})\to\mathcal{B}(\mathcal{K})$ a quantum channel.
The \emph{OPH repair map} is
\[
  \mathcal{R}_{\sigma,\mathcal{N}}(\rho)
  := \sigma^{1/2}\,
     \mathcal{N}^\dagger\!\bigl(
       \mathcal{N}(\sigma)^{-1/2}\,\rho\,\mathcal{N}(\sigma)^{-1/2}
     \bigr)
     \,\sigma^{1/2},
\]
where $\mathcal{N}^\dagger$ is the adjoint channel and inverses are
taken on $\mathrm{supp}(\mathcal{N}(\sigma))$.
\end{definition}

\begin{remark}[Petz map vs.\ trace-distance projection]
The closest-point trace-distance projection
$\mathcal{P}_{\mathcal{S}}(\rho):=\arg\min_{\tau\in\mathcal{S}}\tfrac12\|\rho-\tau\|_1$
is a different object from the Petz map: it is defined by a variational
problem in trace-norm geometry and is not \textrm{CPTP} in general.
The two coincide only in very special cases not automatic in the
\textrm{OPH} setting.
All subsequent properties refer exclusively to
Definition~\ref{def:oph-repair-map}.
\end{remark}

\begin{proposition}[Petz map CPTP: domain-restricted statement
  {\normalfont\small [Established, subject to domain restriction]}]
\label{prop:petz-cptp}
Let $\sigma$ have full support on $\mathcal{H}$.
\begin{enumerate}[(a)]
  \item $\mathcal{R}_{\sigma,\mathcal{N}}$ is completely positive.
  \item $\mathcal{R}_{\sigma,\mathcal{N}}$ is trace-preserving on
        $\mathrm{supp}(\mathcal{N}(\sigma))$, i.e., on inputs $\rho$
        for which $\mathcal{N}(\sigma)^{-1/2}\rho\,\mathcal{N}(\sigma)^{-1/2}$
        is well-defined.
  \item If additionally $\mathcal{N}(\sigma)$ has full rank on $\mathcal{K}$,
        then $\mathcal{R}_{\sigma,\mathcal{N}}$ is \textrm{CPTP} on all
        of $\mathcal{B}(\mathcal{K})$.
\end{enumerate}
If $\mathcal{N}(\sigma)$ is \emph{not} full rank on $\mathcal{K}$,
then either (i)~the domain must be restricted to
$\mathrm{supp}(\mathcal{N}(\sigma))$, or
(ii)~pseudoinverses must replace the inverses (generalised Petz map;
cf.~\cite{Junge2018}), or
(iii)~a regularisation
$\mathcal{N}(\sigma)\mapsto\mathcal{N}(\sigma)+\varepsilon\mathbf{1}$
must be introduced.
Note that full-rank $\sigma$ does \emph{not} prevent $\mathcal{N}(\sigma)$
from being rank-deficient: the channel may map the support of $\sigma$
into a strict subspace of $\mathcal{K}$.
In the \textrm{OPH} setting, whether $\mathcal{N}(\sigma)$ is full rank
depends on the specific overlap channel and must be verified for the chosen analytic channel
model. The finite OPH repair theorem instead uses the declared Lyapunov descent law on a finite
patch net.
\end{proposition}

\begin{proof}
Complete positivity follows from composing three \textrm{CP} operations:
\begin{enumerate}[(i)]
  \item sandwiching by
        $\mathcal{N}(\sigma)^{-1/2}(\cdot)\mathcal{N}(\sigma)^{-1/2}$
        on $\mathrm{supp}(\mathcal{N}(\sigma))$;
  \item $\mathcal{N}^\dagger$;
  \item sandwiching by $\sigma^{1/2}(\cdot)\sigma^{1/2}$.
\end{enumerate}
Trace preservation in the full-rank case:
Petz~\cite{Petz1986}; Fagnola--Umanit\`{a}~\cite{FagnolaUmanita2010}.
\end{proof}

\begin{proposition}[Analytic contraction certificate {\normalfont\small [Assumption-dependent]}]
\label{prop:petz-contraction}
Suppose a declared quotient repair map \(T:Q\to Q\) on a metric physical quotient
\((Q,d_Q)\) is strictly contractive with coefficient \(\lambda\in(0,1)\):
\[
  d_Q(Tx,Ty)\le \lambda d_Q(x,y).
\]
Then \(T\) has at most one fixed point. If a fixed point \(x_\star\) exists, the ideal iterates
obey \(d_Q(T^t x,x_\star)\le \lambda^t d_Q(x,x_\star)\).
This is an analytic contraction condition. It is not required for the finite OPH normal-form
theorem, where termination follows from strict Lyapunov descent of accepted repairs.
\end{proposition}

\begin{theorem}[Noisy approximate repair stability {\normalfont\small [Contraction branch]}]
\label{thm:noisy-contraction}
Assume the contraction certificate of Proposition~\ref{prop:petz-contraction}, and let
\(\widetilde T:Q\to Q\) be an implemented noisy repair map satisfying
\[
d_Q(\widetilde T x,Tx)\le \varepsilon
\qquad\text{for all }x\in Q.
\]
If \(x_\star\) is the ideal fixed point of \(T\), then
\[
d_Q(\widetilde T^t x,x_\star)
\le
\lambda^t d_Q(x,x_\star)+\frac{\varepsilon}{1-\lambda}.
\]
Thus the noisy branch converges only to a controlled error ball, and only after the contraction
certificate and uniform implementation-error bound are supplied.
\end{theorem}

\begin{proof}
The recursion
\[
d_Q(\widetilde T x,x_\star)
\le d_Q(\widetilde T x,Tx)+d_Q(Tx,Tx_\star)
\le \varepsilon+\lambda d_Q(x,x_\star)
\]
iterates to the displayed geometric-series bound.
\end{proof}

\begin{proposition}[Spectral-gap criterion {\normalfont\small [Model-dependent]}]
\label{conj:spectral-gap}
Let \(\mathcal T\) be the Markov, channel, or transfer operator induced by iterated OPH repair on
a declared analytic realization. If \(\mathcal T\) has stationary projection \(\Pi_\star\) and
constants \(C<\infty\), \(\delta>0\) such that on the nonstationary subspace
\[
\|\mathcal T^t-\Pi_\star\|\le C e^{-\delta t},
\]
then the corresponding analytic channel model has exponential convergence. The finite OPH repair
package supplies termination by Lyapunov descent on its declared finite state space; a spectral
gap is a separate quantitative mixing condition for this BFT/QECC-style extension.
\end{proposition}

\begin{theorem}[Exponential Convergence
  {\normalfont\small [Under Proposition~\ref{conj:spectral-gap}]}]
\label{thm:exp-convergence}
Under Proposition~\ref{conj:spectral-gap}, for any initial analytic state \(\rho\),
\[
  \bigl\|\mathcal T^t\rho-\Pi_\star\rho\bigr\|
  \leq C\,e^{-\delta t}\bigl\|\rho-\Pi_\star\rho\bigr\|.
\]
This theorem belongs only to the spectral-gap branch. It is not a consequence of a bare finite
overlap graph or of finite Lyapunov descent alone.
\end{theorem}

% ------------------------------------------------------------
\subsection*{B.3\quad Theorem 3: QECC Correspondence}
% ------------------------------------------------------------

\paragraph{Notation.}
$N=\dim(\mathcal{H})=2^n$ for $n$ physical qubits.
Standard notation: $[[n,k,d]]$ stabilizer code; $K=2^k$; quantum
Singleton bound: $k\leq n-2(d-1)$.

\begin{theorem}[No free min-cut theorem for bare overlap graphs
  {\normalfont\small [Established]}]
\label{thm:no-free-mincut}
Let \(G=(V,E)\) be any connected graph with \(|V|\ge2\). The graph \(G\) alone does not determine
the Hamming distance of the consistency set \(C\) of a finite overlap net on \(G\). In particular,
the same graph can realize a binary constraint code of distance \(1\) or a binary repetition code
of distance \(|V|\).
\end{theorem}

\begin{proof}
Set \(S_i=\{0,1\}\) for every vertex. For the repetition realization, choose \(I_e=\{0,1\}\) and
let both endpoint readouts be the identity. Then every edge imposes \(x_i=x_j\), so the only
global codewords are \(00\cdots0\) and \(11\cdots1\), whose Hamming distance is \(|V|\).

For the trivial-overlap realization, keep the same vertex state spaces but let every endpoint
readout be the constant map to \(0\). Then every binary assignment is globally consistent, so the
minimum Hamming distance among distinct codewords is \(1\). The graph is unchanged. Therefore
distance is a property of the code realization--state spaces, readout maps, logical dictionary,
metric, and error model--not of the bare overlap graph.
\end{proof}

\begin{definition}[Topological-code realization certificate]\label{def:tcert}
An OPH overlap network may be treated as a QECC/topological code only after supplying a tuple
\[
\mathsf{TCert}=
(K,\mathcal H_{\mathrm{phys}},\mathcal H_{\mathrm{code}},
\partial_2,\partial_1,S_X,S_Z,\mathcal L_X,\mathcal L_Z,\mathcal E,\mathcal R),
\]
where \(C_2\xrightarrow{\partial_2}C_1\xrightarrow{\partial_1}C_0\) is a chain complex over
\(\mathbb F_2\), the physical carriers live in \(\mathcal H_{\mathrm{phys}}\), the protected
subspace is \(\mathcal H_{\mathrm{code}}\), \(S_X,S_Z\) are stabilizer or gauge checks,
\(\mathcal L_X,\mathcal L_Z\) are logical-operator classes, \(\mathcal E\) is a declared error
family, and \(\mathcal R\) is a recovery map or recovery family.
\end{definition}

\begin{theorem}[Certified topological-code distance and min-cut
  {\normalfont\small [Assumption-dependent]}]
\label{thm:tcert-distance}
Suppose a certificate \(\mathsf{TCert}\) of Definition~\ref{def:tcert} is supplied, with logical
classes identified as
\[
\mathcal L_X\simeq H_1(K;\mathbb F_2),
\qquad
\mathcal L_Z\simeq H^1(K;\mathbb F_2),
\]
and with boundary conditions excluding lower-weight trivial representatives. Then the certified
distance is
\[
d=
\min\left\{
\min_{\ell\in\mathcal L_X\setminus0}|\ell|,
\min_{\ell^\star\in\mathcal L_Z\setminus0}|\ell^\star|
\right\}.
\]
Only in geometries where this homological systole equals the relevant graph min-cut may one write
\(d=\mathrm{mincut}(G_{\mathrm{OPH}})\) \cite{Kitaev2003,Dennis2002}.
\end{theorem}

\begin{proof}
This is the standard stabilizer/topological-code distance statement once the chain complex,
checks, logical representatives, and boundary conditions are declared. The min-cut equality is an
additional geometric identification of that homological minimum with a graph cut. By
Theorem~\ref{thm:no-free-mincut}, it cannot be inferred from the bare graph.
\end{proof}

\begin{conjecture}[Communication complexity {\normalfont\small [Conjecture]}]
\label{conj:comm-complexity}
The \textrm{OPH} consensus-repair protocol, realised as a quantum
communication task, has per-round complexity $O(n\cdot\mathrm{poly}(d))$
for a chosen communication encoding (cf.~\cite{BCW98}). The fixed finite repair theorem gives
termination after a supplied descent law; it does not by itself fix a quantum communication
complexity class for every implementation.
\end{conjecture}

\begin{theorem}[QECC resilience under a supplied code certificate
  {\normalfont\small [Assumption-dependent]}]
\label{thm:qecc}
Assume a genuine code subspace \(\mathcal H_{\mathrm{code}}\subseteq\mathcal H_{\mathrm{phys}}\)
with projector \(\Pi\), and let \(\mathcal E_t\) be the declared set of errors supported on fewer
than \(t\) corrupted patches or physical carriers. If
\[
\Pi E_a^\dagger E_b \Pi=\alpha_{ab}\Pi
\qquad
\forall E_a,E_b\in\mathcal E_t,
\]
then there exists a recovery channel correcting all errors in \(\mathcal E_t\)
\cite{KL97}. If the supplied certificate also gives distance \(d\), then all errors of weight
\(t<d/2\) are correctable. No such resilience statement follows from the bare overlap graph.
\end{theorem}

\begin{corollary}[QECC extension inventory]\label{cor:qecc-inventory}
Under Definition~\ref{def:tcert}, Theorem~\ref{thm:tcert-distance}, and
Theorem~\ref{thm:qecc}, the OPH BFT/QECC extension carries the following claim split:
\begin{enumerate}[(i)]
  \item \textrm{[Assumption-dependent]} Code distance is the certified homological minimum; it
        equals \(\mathrm{mincut}(G_{\mathrm{OPH}})\) only under the additional systole/min-cut
        identification.
  \item \textrm{[Established]} The Knill--Laflamme \textrm{QECC}
        theorem supplies recovery once the projector and error family satisfy the displayed
        condition.
  \item \textrm{[Conjecture]} Per-round communication complexity is
        $O(n\cdot\mathrm{poly}(d))$.
\end{enumerate}
\end{corollary}

% ------------------------------------------------------------
\subsection*{B.4\quad Theorem 4: Asynchronous Convergence}
% ------------------------------------------------------------

\paragraph{Why fairness alone does not give a probability-1, spectral, or wall-clock statement.}
Standard strong fairness guarantees that every enabled action fires infinitely often along any fair
schedule; it does not impose a probability space on schedules, a transfer-operator spectral gap,
or a message-delay bound. A convergence statement of the form ``converges with probability~1''
requires a measure on schedules. Exponential convergence requires the spectral-gap certificate of
Theorem~\ref{thm:exp-convergence}. Bounded wall-clock liveness requires partial synchrony and
quorum assumptions. The \textrm{FLP} impossibility result~\cite{FLP85} confirms that fairness is
insufficient for bounded-time consensus in a fully asynchronous system.

\paragraph{Additional assumptions for a quantitative bound.}
\begin{enumerate}[(B1)]
  \item Finite known bound $\Delta$ on message delay after \textrm{GST}.
  \item Finite bound $\Phi$ on processing rates.
  \item $f < n/3$.
\end{enumerate}

\begin{theorem}[Eventual finite repair termination
  {\normalfont\small [Finite-descent branch]}]
\label{thm:eventual-convergence}
For a finite OPH patch net with a strict Lyapunov-decreasing accepted repair relation, every
maximal repair run terminates after finitely many accepted repairs. The step count is bounded by
the finite value-set bound of Proposition~\ref{prop:finite-step-bound}. If the local-diamond and
repair-completeness clauses also hold, Newman's lemma upgrades this termination statement to the
unique schedule-independent quotient normal form of Theorem~\ref{thm:confluence}.

This theorem is finite descent convergence in repair steps. It is not trace-norm convergence of a
Petz channel, not probability-one convergence over random schedules, not exponential convergence,
and not a wall-clock liveness theorem.
\end{theorem}

\begin{theorem}[Quantitative Convergence
  {\normalfont\small [Assuming (B1)--(B3)]}]
\label{thm:quantitative-convergence}
In a partially synchronous \textrm{OPH} observer network satisfying
(B1)--(B3), after \textrm{GST} every nonfaulty observer reaches consensus
within $T=O(f\cdot\Delta)$ wall-clock time
(by applying the \textrm{DLS} framework~\cite{DLS88}, Thm.~4.4,
to the \textrm{OPH} repair protocol;
requires (B1) and (B2) explicitly and does not follow from fairness alone).
\end{theorem}

% ------------------------------------------------------------
\subsection*{B.5\quad Extension Boundaries}
% ------------------------------------------------------------

\begin{itemize}
  \item A bare OPH overlap graph is a finite constraint-code presentation only. Its graph does
        not determine code distance, correctable error weight, or min-cut resilience.
  \item Analytic spectral-gap and full-rank estimates for a chosen stochastic or channel-level
        BFT/QECC realization are model-specific refinements. They are separate from the finite
        OPH repair theorem, where the accepted repair law supplies Lyapunov descent directly.
  \item Long-run noisy approximate consensus is available only on the fair-block contraction
        branch of Theorem~\ref{thm:noisy-fair-block-consensus}: the chosen implementation must
        certify fair blocks, expected contraction toward the exact quotient normal-form set, and
        controlled within-block excursions. Fairness alone does not supply that certificate.
  \item Topological-code distance equals graph min-cut only after a concrete chain complex,
        boundary condition, logical-operator dictionary, error family, and systole/min-cut
        identification are supplied for the chosen code realization.
  \item Communication-complexity bounds require a concrete quantum communication encoding and
        implementation cost model. They are not consequences of finite normal-form termination
        alone.
  \item The core OPH consensus paper supplies observable-level confluence and refinement-limit
        normal-form/holonomy classes on their declared theorem surfaces; the BFT/QECC statements
        above are separate protocol-style extensions.
\end{itemize}

% ------------------------------------------------------------
%  Bibliography entries for this appendix
%  (merge into main .bib file or \bibliography)
% ------------------------------------------------------------
%
% @article{FLP85,
%   author  = {Fischer, M.J. and Lynch, N.A. and Paterson, M.S.},
%   title   = {Impossibility of Distributed Consensus with One Faulty Process},
%   journal = {Journal of the ACM},
%   volume  = {32}, number = {2}, pages = {374--382}, year = {1985}}
%
% @article{Lamport1982,
%   author  = {Lamport, L. and Shostak, R. and Pease, M.},
%   title   = {The {Byzantine} Generals Problem},
%   journal = {ACM Transactions on Programming Languages and Systems},
%   volume  = {4}, number = {3}, pages = {382--401}, year = {1982}}
%
% @article{DLS88,
%   author  = {Dwork, C. and Lynch, N.A. and Stockmeyer, L.},
%   title   = {Consensus in the Presence of Partial Synchrony},
%   journal = {Journal of the ACM},
%   volume  = {35}, number = {2}, pages = {288--323}, year = {1988}}
%
% @inproceedings{Castro1999,
%   author    = {Castro, M. and Liskov, B.},
%   title     = {Practical {Byzantine} Fault Tolerance},
%   booktitle = {OSDI 1999}, pages = {173--186}, year = {1999}}
%
% @misc{Moniz2020,
%   author = {Moniz, H.},
%   title  = {The {Istanbul BFT} Consensus Algorithm},
%   eprint = {2002.03613}, year = {2020}}
%
% @misc{Saltini,
%   author = {Saltini, R. and others},
%   title  = {{QBFT} Formal Specification and Verification},
%   howpublished = {\url{https://github.com/Consensys/qbft-formal-spec-and-verification}}}
%
% @article{Petz1986,
%   author  = {Petz, D.},
%   title   = {Sufficient subalgebras and the relative entropy of states
%              of a von Neumann algebra},
%   journal = {Communications in Mathematical Physics},
%   volume  = {105}, number = {1}, pages = {123--131}, year = {1986}}
%
% @article{FagnolaUmanita2010,
%   author  = {Fagnola, F. and Umanit\`{a}, V.},
%   title   = {Generators of detailed balance quantum {Markov} semigroups},
%   journal = {Infinite Dimensional Analysis, Quantum Probability},
%   volume  = {13}, number = {3}, pages = {459--486}, year = {2010}}
%
% @article{Junge2018,
%   author  = {Junge, M. and others},
%   title   = {Universal recovery maps and approximate sufficiency of
%              quantum relative entropies},
%   journal = {Annales Henri Poincar\'e},
%   volume  = {19}, number = {8}, pages = {2505--2555}, year = {2018}}
%
% @article{KL97,
%   author  = {Knill, E. and Laflamme, R.},
%   title   = {Theory of quantum error-correcting codes},
%   journal = {Physical Review A},
%   volume  = {55}, number = {2}, pages = {900--911}, year = {1997}}
%
% @phdthesis{Gottesman1997,
%   author = {Gottesman, D.},
%   title  = {Stabilizer Codes and Quantum Error Correction},
%   school = {California Institute of Technology}, year = {1997}}
%
% @article{Kitaev2003,
%   author  = {Kitaev, A.},
%   title   = {Fault-tolerant quantum computation by anyons},
%   journal = {Annals of Physics},
%   volume  = {303}, number = {1}, pages = {2--30}, year = {2003}}
%
% @article{Dennis2002,
%   author  = {Dennis, E. and Kitaev, A. and Landahl, A. and Preskill, J.},
%   title   = {Topological quantum memory},
%   journal = {Journal of Mathematical Physics},
%   volume  = {43}, number = {9}, pages = {4452--4505}, year = {2002}}
%
% @inproceedings{BCW98,
%   author    = {Buhrman, H. and Cleve, R. and Wigderson, A.},
%   title     = {Quantum vs.\ Classical Communication and Computation},
%   booktitle = {STOC 1998}, pages = {63--68}, year = {1998}}
%
% @book{Attiya2004,
%   author    = {Attiya, H. and Welch, J.},
%   title     = {Distributed Computing: Fundamentals, Simulations, and
%                Advanced Topics},
%   publisher = {Wiley-Interscience}, year = {2004}}
